\section{依恋类型与性：为什么"越靠近越想逃、越爱越焦虑"}

同一句"我想你"，有人听了更亲密，有人听了想后退；同样是短暂的性生活空白，有人能坦然沟通，有人会陷入"他是不是不爱我了"的恐慌。这些差异，很大程度上来自\textbf{依恋风格}——我们在亲密关系中"如何靠近、如何被安抚"的默认模式。

\subsection{四种依恋风格}

\begin{itemize}
	\item \textbf{安全型}：既能亲近，也能独处；能直接表达需求，也能接受对方的"不"。在性关系中更少焦虑，性满意度通常较高。
	\item \textbf{焦虑型}：渴望高度亲密，害怕被抛弃；对伴侣的性兴趣、回复速度、情绪变化高度敏感。性常被用来\textbf{确认被爱}，一旦对方性致不高，容易解读为"不爱我"。
	\item \textbf{回避型}：重视独立，习惯压抑需求、回避冲突；亲密到一定程度会感到被"吞没"，可能用减少性接触、性后回避来恢复心理距离。
	\item \textbf{混乱型（恐惧型）}：既渴望又害怕亲密，常与早年创伤或照料不稳定有关，在性关系中可能忽冷忽热。
\end{itemize}

\begin{tcolorbox}[colback=blue!5!white,colframe=blue!50!black,title=补充说明]
依恋风格不是标签，也不是命运。它是"倾向"，会随关系经历、治疗与自我觉察而改变；同一个人在不同关系中、对不同伴侣，也可能表现出不同风格。
\end{tcolorbox}

\subsection{依恋风格如何在性里"显形"}

\begin{itemize}
	\item \textbf{发起与拒绝}：焦虑型伴侣常是更主动的一方，被拒绝时的痛苦更强；回避型伴侣的"不想做"常被误读为嫌弃，其实可能是需要空间。
	\item \textbf{性后的时刻}：安全型能享受事后亲密；回避型可能在结束后想独处、不想交流，被对方解读为"用完就走"；焦虑型则渴望拥抱与确认，得不到会格外失落。
	\item \textbf{冲突中的性}：有人用"性和好"（和解式性爱）化解矛盾，有人一吵架就"性趣全无"——背后是"靠近"与"撤出"两套策略的碰撞。
\end{itemize}

\subsection{如何相处：把"回避--追逐"的死循环拆开}

\begin{itemize}
	\item \textbf{识别循环而非指责}："你追得越紧，我退得越远；我退得越远，你追得越紧"——这是两人共同制造的循环，不是谁的错。
	\item \textbf{焦虑方}：学习自我安抚（而非立刻向伴侣求证），把"你不做就是不爱我"翻译成"我现在需要一点确认"。
	\item \textbf{回避方}：练习"给一句安心的回应再退开"——如"我现在有点想静静，但我很爱你，晚点我来找你"。这能大幅降低对方的恐慌。
	\item \textbf{安全型的伴侣}可以起到"情绪稳压器"的作用：稳定、可预测的回应，会慢慢让伴侣的依恋模式向安全靠拢。
	\item \textbf{若一方有未处理的创伤}（性侵、暴力、被忽视），依恋问题常与创伤交织，应在专业心理治疗下处理，而非靠"多沟通"硬扛。
\end{itemize}

\begin{tcolorbox}[colback=pink!5!white,colframe=red!50!black,title=贴心小叮咛]
性关系里的很多"作"与"冷"，不是品质问题，而是\textbf{依恋系统在报警}。看懂对方是"需要一点空间"还是"需要一点确认"，比争论"你到底爱不爱我"有用得多。若两人总在同一个循环里打转，寻求伴侣治疗往往比反复争吵更高效。
\end{tcolorbox}
